% Options for packages loaded elsewhere
\PassOptionsToPackage{unicode}{hyperref}
\PassOptionsToPackage{hyphens}{url}
\PassOptionsToPackage{dvipsnames,svgnames,x11names}{xcolor}
\documentclass[
  11pt,
]{article}
\usepackage{xcolor}
\usepackage[margin=27mm]{geometry}
\usepackage{amsmath,amssymb}
\setcounter{secnumdepth}{-\maxdimen} % remove section numbering
\usepackage{iftex}
\ifPDFTeX
  \usepackage[T1]{fontenc}
  \usepackage[utf8]{inputenc}
  \usepackage{textcomp} % provide euro and other symbols
\else % if luatex or xetex
  \usepackage{unicode-math} % this also loads fontspec
  \defaultfontfeatures{Scale=MatchLowercase}
  \defaultfontfeatures[\rmfamily]{Ligatures=TeX,Scale=1}
\fi
\usepackage{lmodern}
\ifPDFTeX\else
  % xetex/luatex font selection
\fi
% Use upquote if available, for straight quotes in verbatim environments
\IfFileExists{upquote.sty}{\usepackage{upquote}}{}
\IfFileExists{microtype.sty}{% use microtype if available
  \usepackage[]{microtype}
  \UseMicrotypeSet[protrusion]{basicmath} % disable protrusion for tt fonts
}{}
\makeatletter
\@ifundefined{KOMAClassName}{% if non-KOMA class
  \IfFileExists{parskip.sty}{%
    \usepackage{parskip}
  }{% else
    \setlength{\parindent}{0pt}
    \setlength{\parskip}{6pt plus 2pt minus 1pt}}
}{% if KOMA class
  \KOMAoptions{parskip=half}}
\makeatother
\usepackage{longtable,booktabs,array}
\newcounter{none} % for unnumbered tables
\usepackage{calc} % for calculating minipage widths
% Correct order of tables after \paragraph or \subparagraph
\usepackage{etoolbox}
\makeatletter
\patchcmd\longtable{\par}{\if@noskipsec\mbox{}\fi\par}{}{}
\makeatother
% Allow footnotes in longtable head/foot
\IfFileExists{footnotehyper.sty}{\usepackage{footnotehyper}}{\usepackage{footnote}}
\makesavenoteenv{longtable}
\setlength{\emergencystretch}{3em} % prevent overfull lines
\providecommand{\tightlist}{%
  \setlength{\itemsep}{0pt}\setlength{\parskip}{0pt}}
\usepackage{bookmark}
\IfFileExists{xurl.sty}{\usepackage{xurl}}{} % add URL line breaks if available
\urlstyle{same}
\hypersetup{
  pdftitle={Rational 7-torsion and the 2-primary torsion of geometrically split Jacobians},
  colorlinks=true,
  linkcolor={Maroon},
  filecolor={Maroon},
  citecolor={Blue},
  urlcolor={Blue},
  pdfcreator={LaTeX via pandoc}}

\title{Rational 7-torsion and the 2-primary torsion of geometrically
split Jacobians}
\author{}
\date{11 November 2026}

\begin{document}
\maketitle

\section{Theorem and scope}\label{theorem-and-scope}

\textbf{Theorem.} Let \(A/\mathbb Q\) be a principally polarized abelian
surface that is simple over \(\mathbb Q\) and geometrically split. If
\(A(\mathbb Q)\) has a point of order \(7\), then \[
\#A(\mathbb Q)[2^\infty]\leq 4.
\]

\textbf{Corollary.} No genus-\(2\) Jacobian in this class has rational
torsion containing \(C_2\times C_{28}\). In particular, the exact group
\([2,28]\) is impossible in the census column \textbf{geometrically
split, simple over \(\mathbb Q\)}. More generally, every group with
\(7\)-torsion and \(2\)-primary subgroup of order at least eight is
excluded in this class, including \([56]\) and \([2,2,14]\).

Here geometrically split means that \(A_{\overline{\mathbb Q}}\) is
isogenous to a product of elliptic curves. There is no restriction on
the splitting degree or the integral endomorphism ring. The principal
polarization is used in the cyclic cubic case. The theorem does not
exclude \([2,28]\) for Jacobians split over \(\mathbb Q\), or for
geometrically simple Jacobians. The combined geometrically split
existence and infinitude question for \([2,28]\) therefore remains open.
The bound also leaves \([28]\) and \([2,14]\) unexcluded.

The complete proof is available as
\href{https://genus-2-torsion.github.io/data/knowledge/sources/qsimple_228_obstruction.pdf}{PDF}
and
\href{https://genus-2-torsion.github.io/data/knowledge/sources/qsimple_228_obstruction.tex}{LaTeX}.
The exact
\href{https://genus-2-torsion.github.io/data/knowledge/sources/qsimple_228_checks.m}{Magma
certificate} and its
\href{https://genus-2-torsion.github.io/data/knowledge/sources/qsimple_228_checks.log}{completed
log} accompany the proof.

\section{1. Conventions and the existing
reduction}\label{conventions-and-the-existing-reduction}

For a number field \(k\), put
\(G_k=\operatorname{Gal}(\overline{\mathbb Q}/k)\). Write \(T_\ell\) for
the Tate module, \(V_\ell=T_\ell\otimes\mathbb Q_\ell\), and
\(\chi_\ell\) for the cyclotomic character. We use residual
semisimplification to compare different stable lattices in one rational
Tate representation; rational torsion is not moved along arbitrary
isogenies.

The
\href{https://genus-2-torsion.github.io/proofs/qsimple-2214-obstruction.html}{published
rational-\(7\)-torsion and \(2\)-rank proof}, especially Proposition 2,
the cyclic-field reduction, and its remarks on \([2,28]\), gives the
following inputs.

\begin{itemize}
\tightlist
\item
  If the geometric elliptic factors are nonisogenous, or form a non-CM
  square with quadratic endomorphism field, rational \(7\)-torsion
  forces the entire rational \(2\)-primary torsion to have order at most
  four. This uses Kenku--Momose's quadratic elliptic torsion theorem.
\item
  For a CM elliptic square, a point count at \(3\) gives the same bound.
\item
  The only remaining endomorphism fields are cyclic of degree \(3\) or
  \(6\).
\end{itemize}

In these last two cases put \[
D=\operatorname{End}^0(A_{\overline{\mathbb Q}})\simeq M_2(\mathbb Q),
\qquad F=\operatorname{End}^0_{\mathbb Q}(A)=\mathbb Q(\sqrt{-3}),
\] and let \(L\) be the smallest field over which all geometric
endomorphisms are defined. The normalization and cubic untwisting below
are those proved in the
\href{https://genus-2-torsion.github.io/proofs/qsimple-63-obstruction.html}{order-\(63\)
proof}, Sections 5--7 and 10.2, and the
\href{https://genus-2-torsion.github.io/proofs/geometric-26-39-obstruction.html}{order-\(26/39\)
proof}, Section 8. We show that, in either remaining case, rational
\(7\)-torsion is incompatible with rational \(2\)-torsion.

\section{2. Cyclic cubic endomorphism field: polarization at
3}\label{cyclic-cubic-endomorphism-field-polarization-at-3}

Suppose \(\operatorname{Gal}(L/\mathbb Q)=C_3\). Cubic untwisting
expresses the rational Tate representation as \[
V_\ell(A)\simeq V_\ell(E_0)\otimes_{\mathbb Q_\ell}W_\ell,
\] where \(E_0/\mathbb Q\) is non-CM and \(W\) is the two-dimensional
rational representation given by a faithful character \[
\psi:G_{\mathbb Q}\longrightarrow\mu_3\subset F^\times,
\qquad \ker\psi=G_L.
\] This is the construction in Section 10.2 of the order-\(63\) proof
before its specialization to an elliptic isogeny class. Replacing
\(\psi\) by its inverse does not affect the argument.

\subsection{\texorpdfstring{2.1. Principal polarization forces a
rational 3-isogeny on
\(E_0\)}{2.1. Principal polarization forces a rational 3-isogeny on E\_0}}\label{principal-polarization-forces-a-rational-3-isogeny-on-e_0}

Suppose that \(E_0[3]\) is irreducible over \(\mathbb F_3\). Let \[
\mathcal O=\mathbb Z_3[\zeta_3],\qquad \pi=\zeta_3-1,
\qquad T_0=T_3(E_0)\otimes_{\mathbb Z_3}\mathcal O.
\] The lattice \(T_0\) is \(G_{\mathbb Q}\)-stable, with
\(G_{\mathbb Q}\) acting on its second factor by \(\psi\). Since
\(\pi^2\) is a unit times \(3\), the module \(T_0/3T_0\) is an extension
of \(W_0=E_0[3]\) by itself, with natural submodule \(\pi T_0/3T_0\).

This extension is nonsplit. Modulo \(\pi^2\), write \[
\psi(g)=1+c(g)\pi,
\] where \(c:G_{\mathbb Q}\to\mathbb F_3\) is a nonzero additive
character. A \(G_{\mathbb Q}\)-linear splitting would give an
\(\mathbb F_3\)-linear endomorphism \(S\) of \(W_0\) satisfying, up to
an irrelevant overall sign, \[
\rho(g)^{-1}S\rho(g)-S=c(g)I.
\] Taking traces gives \(2c(g)=0\) in \(\mathbb F_3\), a contradiction.
As \(W_0\) is irreducible, a nonsplit length-two self-extension has a
unique proper nonzero submodule.

It follows that every \(G_{\mathbb Q}\)-stable \(\mathbb Z_3\)-lattice
in this rational representation is \(\pi^jT_0\) for some integer \(j\).
Indeed, scale such a lattice \(T\) by a power of \(3\) so that
\(T\subseteq T_0\) but \(T\nsubseteq3T_0\). If its image modulo \(3T_0\)
is the whole module, Nakayama gives \(T=T_0\). Otherwise its image is
the unique proper nonzero submodule, so \(T\subseteq\pi T_0\). Then
\(\pi^{-1}T\subseteq T_0\) but \(\pi^{-1}T\nsubseteq\pi T_0\). Its
reduction cannot be that proper submodule, and Nakayama gives
\(\pi^{-1}T=T_0\).

The space of \(G_{\mathbb Q}\)-equivariant alternating forms on
\(V_3(A)\) with cyclotomic similitude is one-dimensional over
\(\mathbb Q_3\). To see this, restrict first to \(G_L\). By Faltings's
theorem and \(\operatorname{End}_L(E_0)=\mathbb Z\), equivariant
bilinear forms on \(V_3(E_0)\) with cyclotomic multiplier form the line
spanned by the elliptic Weil pairing \(e\). Thus the alternating forms
on \(V_3(A)\) have the form \(e\otimes b\), with \(b\) symmetric on the
multiplicity space. Invariance under \(\mu_3\) forces \(b\) to be a
scalar multiple of \[
b_0(a,b)=\operatorname{Tr}_{\mathbb Q_3(\zeta_3)/\mathbb Q_3}(a\overline b).
\]

For this trace pairing, the dual of \(\mathcal O\) is the inverse
different \(\pi^{-1}\mathcal O\). Since \(T_3(E_0)\) is self-dual for
\(e\), the dual of \(\pi^jT_0\) under \(q(e\otimes b_0)\), for
\(q\in\mathbb Q_3^\times\), is \[
\pi^{-j-1-2v_3(q)}T_0.
\] Self-duality would require \[
2j=-1-2v_3(q),
\] which is impossible by parity. A principal polarization supplies
exactly such a perfect alternating form on \(T_3(A)\). Hence \(E_0[3]\)
is reducible, and \(E_0\) has a rational cyclic \(3\)-isogeny. The
argument in fact excludes every polarization of degree prime to \(3\) in
the isogeny class under the irreducibility hypothesis.

\subsection{\texorpdfstring{2.2. Rational 7-torsion and
\(X_0(21)\)}{2.2. Rational 7-torsion and X\_0(21)}}\label{rational-7-torsion-and-x_021}

Since \(7\) splits in \(F\), the two characters of \(W\) over
\(\mathbb F_7\) are \(\psi\) and \(\psi^{-1}\). A trivial constituent
coming from \(A(\mathbb Q)[7]\ne0\) forces \(E_0[7]\) to have one of
these character constituents. Thus \(E_0[7]\) is reducible and \(E_0\)
has a rational cyclic \(7\)-isogeny. The coprime kernels of the \(3\)-
and \(7\)-isogenies give a rational cyclic \(21\)-isogeny.

The exact certificate enumerates all rational points of \(X_0(21)\),
using algebraic rank zero and its torsion group \(C_2\times C_4\). Its
four noncuspidal \(j\)-invariants are \[
-\frac{189613868625}{128},\qquad \frac{3375}{2},\qquad
-\frac{1159088625}{2097152},\qquad -\frac{140625}{8}.
\] For each, the elliptic \(2\)-division polynomial is irreducible with
nonsquare discriminant, so its splitting field has Galois group \(S_3\).
These properties are unchanged by quadratic twists; none of the four
values is \(0\) or \(1728\).

An \(S_3\) extension has no cyclic cubic subextension, so it is linearly
disjoint from \(L\). Therefore \(E_0[2]\) remains irreducible on
\(G_L\). The representation \(W\) becomes trivial there, so the residual
semisimplification of \(A[2]\) restricted to \(G_L\) has two copies of
this nontrivial irreducible constituent and has no trivial constituent.
In particular \(A(\mathbb Q)[2]=0\). This completes the cubic case.

\section{3. Cyclic sextic field: a level-21 point with square
discriminant}\label{cyclic-sextic-field-a-level-21-point-with-square-discriminant}

Suppose \(\operatorname{Gal}(L/\mathbb Q)=C_6\), and let \(K\) be its
quadratic subfield. The construction in Section 8 of the order-\(26/39\)
proof gives the following facts without a torsion assumption.

\begin{enumerate}
\def\labelenumi{\arabic{enumi}.}
\tightlist
\item
  Both \(L\) and \(K\) are totally real.
\item
  After an isogeny normalizing the action of \(\mathbb Z[\zeta_3]\) and
  a cubic twist, one obtains \(B/\mathbb Q\) whose geometric
  endomorphisms are all defined over \(K\).
\item
  There is \(\theta\in\operatorname{End}_{\mathbb Q}(B)\) with
  \(\theta^2=-3\). For each elliptic \(E_0\subset B_K\), the restriction
  of \(\theta\) gives a cyclic \(K\)-isogeny of degree \(3\) to its
  distinct conjugate \(E_0^\tau\).
\item
  On \(G_K\), the original representation differs from that of \(B\) by
  a faithful cubic multiplicity character \(\psi\), with kernel \(G_L\).
\end{enumerate}

\subsection{3.1. The level-7 structure and its
conjugation}\label{the-level-7-structure-and-its-conjugation}

Rational \(7\)-torsion supplies a one-dimensional constituent in one of
the two \(7\)-adic \(F\)-blocks. The lattice argument in Proposition 9.1
of the order-\(26/39\) proof applies with \(7\) in place of \(13\). Here
are the steps needed for this application.

\begin{itemize}
\tightlist
\item
  In a chosen \(F\)-block of the normalized surface, cubic twisting
  preserves reducibility modulo \(7\). Choose a \(G_{\mathbb Q}\)-stable
  line in that residual block of \(B\).
\item
  Inclusion of \(E_0\) followed by projection to the chosen block is a
  \(G_K\)-equivariant isomorphism on rational Tate modules. Faltings's
  theorem gives irreducibility of \(V_7(E_0)\), and the projection
  cannot vanish: otherwise \(V_7(E_0)\) would be stable under
  \(\theta\), contradicting \(\theta E_0\ne E_0\).
\item
  Realize the block lattice by a \(K\)-isogeny \(h:E_0\to E\) of degree
  \(7^b\) and pull the residual line back to a \(G_K\)-stable subgroup
  \(C_7\subset E[7]\).
\item
  Galois action outside \(G_K\) gives a \(G_K\)-equivariant isomorphism
  \(T_7(E^\tau)\to T_7(E)\). By Faltings's theorem it is a
  \(\mathbb Z_7\)-unit times a primitive generator of
  \(\operatorname{Hom}_K(E^\tau,E)\), whose degree is prime to \(7\).
\item
  Transporting the degree-three isogeny through \(h\) gives an isogeny
  of degree \(3\cdot7^{2b}\). Its primitive generator must consequently
  have degree \(3\). The Galois-stable residual line is respected by
  this generator.
\end{itemize}

Reversing the direction if necessary, we obtain a cyclic isogeny
\(\nu:E\to E^\tau\) of degree \(3\) with \(\nu(C_7)=\tau C_7\). Its
kernel and \(C_7\) define \[
x=(E,\ker\nu+C_7)\in Y_0(21)(K),\qquad \tau(x)=w_3(x).
\] The last equality is the moduli definition of \(w_3\): quotient by
\(\ker\nu\) sends \(E[3]/\ker\nu\) to the kernel of the conjugate
degree-three isogeny, and sends \(C_7\) to \(\tau C_7\). This is the
same compatibility argument used at level \(39\). The point \(x\) cannot
be rational, since a \(w_3\)-fixed point would give a degree-three
endomorphism of a non-CM elliptic curve.

\subsection{3.2. The elliptic 2-division
field}\label{the-elliptic-2-division-field}

Suppose \(A(\mathbb Q)[2]\ne0\). After extension of scalars to
\(\mathbb F_4\), the multiplicity characters on \(G_K\) are \(\psi\) and
\(\psi^{-1}\). A trivial residual constituent forces
\(E[2]\otimes\mathbb F_4\) to have one of these nontrivial cubic
characters as a constituent. Its determinant over \(\mathbb F_2\) is
\(1\), so its two characters are \(\psi\) and \(\psi^{-1}\). The only
subgroup of \(\operatorname{GL}_2(\mathbb F_2)=S_3\) with these
constituents is \(C_3\); the natural \(S_3\) module remains irreducible
over \(\mathbb F_4\). The intervening \(7\)-power isogeny leaves
\(E[2]\) unchanged. Thus the \(2\)-division field of \(E\) is exactly
\(L/K\), and its discriminant is a square in \(K\).

For \(j(E)\notin\{0,1728\}\) the identity \[
j(E)-1728=\frac{c_6(E)^2}{\Delta(E)}
\] shows that \(j(E)-1728\) is a square in \(K\). Hence \(x\) lifts to a
\(K\)-point of the double cover of \(X_0(21)\) obtained by adjoining
\(\sqrt{j-1728}\). The next section determines the relevant points of
this genus-three cover.

\section{4. The discriminant cover and its rank-zero
quotients}\label{the-discriminant-cover-and-its-rank-zero-quotients}

Use the standard model \[
X_0(21):y^2+xy=x^3-4x-1,
\qquad t=\frac{y-x-3}{x+1}.
\] The certificate checks that the \(j\)-map has degree \(32\), that its
four cusp pole orders are \(21,7,3,1\), and that the designated \(w_3\)
satisfies \(\Phi_3(j,j\circ w_3)=0\), whereas \(w_7\) and \(w_{21}\) do
not. It also checks the point counts against the Hecke traces at primes
through \(43\) away from \(3\) and \(7\).

In the function field one has \[
x^2-(t^2+3t+3)x-t^2-6t-10=0,
\qquad j-1728=H\cdot(\text{a rational function})^2,
\qquad H=t(x+2).
\] Write the cover as \(z^2=H\). The involution \(w_3\) fixes \(t\). Put
\(q=t/(x+2)\). Then \[
\frac{w_3(H)}{H}=q^2,\qquad q\,w_3(q)=1.
\] The two involutions lifting \(w_3\) have the elliptic quotients \[
\begin{aligned}
C_+:&\quad U^2=t(t^2+5t+7),& U&=z(1+q),\\
C_-:&\quad U^2=t(t^2+t+7),& U&=z(1-q).
\end{aligned}
\] The identities defining these quotients are checked exactly. Both
elliptic curves have rank zero by algebraic higher \(2\)-isogeny
descent. Their full rational torsion groups and the corresponding finite
\(t\)-values are

{\def\LTcaptype{none} % do not increment counter
\begin{longtable}[]{@{}lll@{}}
\toprule\noalign{}
Curve & Rational torsion & Finite \(t\)-values \\
\midrule\noalign{}
\endhead
\bottomrule\noalign{}
\endlastfoot
\(C_+\) & \(C_2\) & \(0\) \\
\(C_-\) & \(C_6\) & \(0,1,7\) \\
\end{longtable}
}

Each also has the point at infinity. The fibers \(t=0,\infty\) on
\(X_0(21)\) are exactly its four cusps. At the remaining values the two
\(x\)-coordinates are \[
\begin{array}{c|c}
t&x\\\hline
1&(7\pm3\sqrt{13})/2\\
7&(73\pm21\sqrt{13})/2.
\end{array}
\]

Since \(\tau(x)=w_3(x)\), conjugation of any lift to the cover agrees
with one of the two lifts of \(w_3\). Its image on the corresponding
elliptic quotient is rational. The complete rational point sets above
therefore force \[
K=\mathbb Q(\sqrt{13}),\qquad t\in\{1,7\}.
\] The corresponding \(j\)-invariants are \[
\begin{aligned}
t=1:&\quad6565897584689283072\pm1821052339310297088\sqrt{13},\\
t=7:&\quad75497472\pm20938752\sqrt{13}.
\end{aligned}
\] They are neither \(0\) nor \(1728\). Thus every elliptic curve over
\(K\) with one of these values is a quadratic twist of the model
\(E_j=\operatorname{EllipticCurveFromjInvariant}(j)\).

\section{5. Frobenius at 29}\label{frobenius-at-29}

The prime \(29\) splits in \(\mathbb Q(\sqrt{13})\). For both displayed
\(j\)-invariants and both primes above \(29\), the certificate checks
that \(E_j\) has good reduction, its \(2\)-division polynomial splits
completely, and its Frobenius trace has absolute value \(6\).

The actual \(E\) is a quadratic twist of \(E_j\), with unchanged
\(2\)-division field. Therefore \(\psi\) is unramified and trivial on
Frobenius at these primes: it is the faithful cubic \(2\)-division
character. If the quadratic twist were ramified, inertia would act on
the original mod-\(7\) representation as \(-I\), since \(E_j\) and
\(\psi\) are both unramified. This contradicts the trivial constituent
supplied by rational \(7\)-torsion. Hence the quadratic twist is
unramified and its trace is \(a\in\{6,-6\}\).

A rational Frobenius at \(29\) belongs to \(G_K\), where the cubic
multiplicity twist is now trivial. The Frobenius polynomial of \(A\) is
consequently \[
P_{29}(T)=(T^2-aT+29)^2.
\] The same representation is unramified, so \(A\) has good reduction at
\(29\). Prime-to-\(29\) rational torsion injects into its reduction, and
its order must divide \[
P_{29}(1)=(30-a)^2\in\{576,1296\}.
\] Neither integer is divisible by \(7\). This contradiction completes
the sextic case.

\section{6. Conclusion}\label{conclusion}

In both the cubic and sextic cases, rational \(7\)-torsion precludes
rational \(2\)-torsion. The entire rational \(2\)-primary group is then
trivial, because every nontrivial finite \(2\)-group has an element of
order two. In all other geometric cases the existing argument bounds its
order by four. This proves the theorem.

The exclusion of \([2,28]\) concerns only the Q-simple geometrically
split column. Its geometrically simple example is retained, its
split-over-\(\mathbb Q\) case is not settled here, and the combined
geometrically split infinitude grade remains open. The same general
bound excludes groups outside the census, such as cyclic \([56]\), in
the Q-simple geometrically split class.

\section{7. Verification}\label{verification}

The standalone certificate
\href{https://genus-2-torsion.github.io/data/knowledge/sources/qsimple_228_checks.m}{qsimple\_228\_checks.m}
checks every new numerical input with an assertion: the modular model
and \(j\)-map, cusps and Hecke traces, identification of \(w_3\), the
genus of the discriminant cover, the exact quotient identities,
algebraic rank-zero upper bounds, complete elliptic torsion sets, the
four rational \(j\)-classes and their \(S_3\) \(2\)-division groups, and
both embeddings above \(29\) for each remaining quadratic fiber. Its
\href{https://genus-2-torsion.github.io/data/knowledge/sources/qsimple_228_checks.log}{log}
ends with \texttt{VERIFIED\_228\_MODULAR\_CERTIFICATE\_DONE}.

The proof has been independently reviewed, and the final certificate was
independently rerun. The finite-field calculations at \(29\) were also
recomputed separately. The completeness statements use algebraic descent
and complete torsion enumeration, with no analytic-rank or BSD
assumption. No bounded search is used as a nonexistence proof.

\section{References}\label{references}

\begin{enumerate}
\def\labelenumi{\arabic{enumi}.}
\tightlist
\item
  \href{https://genus-2-torsion.github.io/proofs/qsimple-2214-obstruction.html}{Rational
  7-torsion and the 2-rank of geometrically split abelian surfaces},
  census proof of 10 November 2026: Proposition 2, the cyclic-field
  reduction, and the remarks on \([2,28]\).
\item
  \href{https://genus-2-torsion.github.io/proofs/qsimple-63-obstruction.html}{No
  rational 63-torsion on a geometrically split abelian surface simple
  over Q}, census proof of 21 October 2026: Sections 5--7 and 10.2,
  including the general cubic untwisting construction.
\item
  \href{https://genus-2-torsion.github.io/proofs/geometric-26-39-obstruction.html}{No
  rational points of order 26 or 39 on geometrically split abelian
  surfaces}, census proof of 21 October 2026: Section 8 and the lattice
  and level-structure steps of Proposition 9.1. The present proof
  repeats the transfer at \(7\) and supplies the new modular-cover
  calculation.
\item
  G. Faltings, \emph{Endlichkeitssätze für abelsche Varietäten über
  Zahlkörpern}, Invent. Math. \textbf{73} (1983), 349--366.
  \href{https://doi.org/10.1007/BF01388432}{Published paper}. We use the
  isogeny and Tate-module endomorphism theorems; in Section 2.1 the
  required consequence is
  \(\operatorname{End}_{G_L}(V_3(E_0))=\mathbb Q_3\) for non-CM \(E_0\).
\item
  J.-P. Serre and J. Tate, \emph{Good reduction of abelian varieties},
  Ann. of Math. \textbf{88} (1968), 492--517.
  \href{https://doi.org/10.2307/1970722}{Primary paper}: the
  good-reduction criterion and prime-to-residue-characteristic torsion
  injection.
\item
  Magma's
  \href{https://magma.maths.usyd.edu.au/magma/handbook/modular_curves}{small
  modular curve and j-map tools}. The model and involution conventions
  are checked in the accompanying script. Its
  \texttt{TwoPowerIsogenyDescentRankBound} computations are algebraic
  upper bounds.
\end{enumerate}

\end{document}
